%!TEX root = ../main.tex
% =============================================================
%  CAPÍTULO 6 — ANÁLISIS DE RIESGOS     (Entrega 6 · 5 % · 29/9)
% =============================================================

\chapter{Análisis de riesgos}\label{cap:riesgos}

\begin{guia}[Guía de la cátedra: qué espera ver en este capítulo]
Formato obligatorio de cada riesgo: \textbf{``Debido a [causa], podría
ocurrir [evento], lo que provocaría [efecto]''}. Sin causa es un miedo,
no un riesgo.
\begin{enumerate}
    \item \textbf{Tabla de riesgos} con ID, descripción causa--evento--efecto,
          naturaleza, etapa, probabilidad (1--5), impacto (1--5),
          \gls{exposicion} (1--25), respuesta, responsable y contingencia;
          ordenada por exposición descendente. Con los 6 a 8 primeros bien
          tratados alcanza; el resto se monitorea.
    \item Al menos un riesgo de \textbf{cada naturaleza} (alcance, técnicos,
          recursos, terceros, datos y cumplimiento) y de \textbf{cada etapa}
          (antes de arrancar, durante, pasaje a producción, operación).
    \item Al menos un riesgo \textbf{externo} (``si no tienen ninguno, no
          miraron para afuera'') y al menos uno con respuesta
          \textbf{aceptar} bien justificado (``aceptar se escribe; ignorar
          se paga'').
    \item \textbf{Supuestos y restricciones} separados de los riesgos. Cada
          supuesto que falla se convierte en riesgo.
    \item Que sean propios: ``un riesgo copiado de internet se nota a tres
          metros''.
\end{enumerate}
Los siete que matan proyectos: alcance que nunca cierra (el número uno
de la materia), pérdida de la fuente de datos, equipo que se desarma,
tecnología que no da, inviabilidad de costo, barrera legal o ética,
nadie a cargo. Y la pregunta de jurado: ¿quién lo mantiene cuando
ustedes se reciben? Marco de referencia: \citet{pmi2021}.
\end{guia}

\section{Metodología}

Este capítulo identifica, analiza y trata los riesgos del proyecto
MotoSafe GBA siguiendo el proceso visto en la cátedra: identificar,
analizar, evaluar y tratar. Un riesgo, a diferencia de un problema,
todavía no ocurrió: es algo que puede llegar a pasar, con una
probabilidad y una consecuencia asociadas. Cuando el evento ya sucedió y
ya se resolvió, como varios de los casos de este capítulo, se lo
documenta igual, porque la forma en que se trató muestra mejor la
capacidad de gestión del proyecto que un riesgo puramente hipotético.

Cada riesgo se redacta explicando su causa, el evento que podría ocurrir
y su efecto (por ejemplo: "debido a X, podría pasar Y, lo que provocaría
Z"), en vez de un enunciado vago del tipo "puede fallar el servidor". La
probabilidad y el impacto se califican en una escala de 1 a 5, y el
producto de ambos (la exposición) ordena la tabla de mayor a menor
prioridad. Cada riesgo recibe una de cuatro respuestas: evitarlo,
reducirlo, transferirlo a un tercero o aceptarlo; aceptar un riesgo no
es lo mismo que ignorarlo, sino dejarlo declarado por escrito con el
motivo.

\begin{longtable}[]{@{}
  >{\raggedright\arraybackslash}p{(\columnwidth - 4\tabcolsep) * \real{0.0997}}
  >{\raggedright\arraybackslash}p{(\columnwidth - 4\tabcolsep) * \real{0.3878}}
  >{\raggedright\arraybackslash}p{(\columnwidth - 4\tabcolsep) * \real{0.5125}}@{}}
\caption{escala de probabilidad e impacto usada en este capítulo (1 a 5). La exposición es su producto, de 1 a 25.}\label{tab:6-0}\\
\toprule\noalign{}
\begin{minipage}[b]{\linewidth}\raggedright
\textbf{Valor}
\end{minipage} & \begin{minipage}[b]{\linewidth}\raggedright
\textbf{Probabilidad}
\end{minipage} & \begin{minipage}[b]{\linewidth}\raggedright
\textbf{Impacto}
\end{minipage} \\
\midrule\noalign{}
\endfirsthead
\toprule\noalign{}
\begin{minipage}[b]{\linewidth}\raggedright
\textbf{Valor}
\end{minipage} & \begin{minipage}[b]{\linewidth}\raggedright
\textbf{Probabilidad}
\end{minipage} & \begin{minipage}[b]{\linewidth}\raggedright
\textbf{Impacto}
\end{minipage} \\
\midrule\noalign{}
\endhead
\bottomrule\noalign{}
\endlastfoot
1 & Muy baja: prácticamente no va a pasar. & Insignificante: no cambia
nada relevante del proyecto. \\
2 & Baja: podría pasar, pero no es lo esperable. & Menor: molesta pero
no afecta ningún resultado clave. \\
3 & Media: tan probable como no. & Moderado: afecta un resultado, pero
se puede compensar. \\
4 & Alta: es lo más probable que pase. & Grave: afecta un resultado
central del proyecto. \\
5 & Muy alta: ya pasó, o va a pasar casi seguro. & Crítico: pone en duda
el aporte principal del proyecto. \\
\end{longtable}

\section{Supuestos y restricciones}

Los supuestos son hechos o hipótesis sobre los que se apoya el proyecto;
las restricciones son condiciones dadas que lo limitan en tiempo, costo
o recursos. Se listan por separado de los riesgos porque no son lo
mismo: si un supuesto resulta falso, invalida parte del análisis,
mientras que un riesgo es un evento que puede o no llegar a ocurrir.

\subsection{Supuestos}

\begin{itemize}
\item
  Los datos oficiales de siniestros de CABA usados para entrenar el
  modelo son representativos del riesgo vial actual, aunque no incluyan
  siniestros posteriores a la última actualización del dataset.
\item
  El usuario final de la aplicación cuenta con un teléfono con GPS y
  conexión de datos al usar el modo de navegación en la vía pública.
\item
  Los servicios externos de los que depende el sistema (ruteo con
  tráfico, clima, túnel de acceso público) sostienen durante el
  desarrollo y la defensa de la tesis el mismo nivel de servicio con el
  que se los integró.
\item
  Los cálculos económicos del Capítulo 5 (cuántos usuarios se suman por
  año, el horizonte de cinco años, la estructura de costos) son una
  proyección de referencia: todavía no hay ninguna validación real de
  mercado que confirme que esos números se van a cumplir.
\end{itemize}

\subsection{Restricciones}

\begin{itemize}
\item
  Restricción de tiempo: el proyecto tiene que entregarse dentro del
  cronograma fijo de diez entregas de la cátedra, lo que limita cuánto
  se puede seguir profundizando cada punto pendiente antes de la entrega
  final.
\item
  Restricción de recursos humanos: el desarrollo está a cargo de una
  sola persona, con asistencia de herramientas de inteligencia
  artificial, lo que limita cuánto se puede paralelizar y revisar de
  forma cruzada.
\item
  Restricción de datos: solo la Ciudad de Buenos Aires publica un
  registro de siniestros con ubicación exacta; el resto del conurbano no
  cuenta con un equivalente accesible.
\item
  Restricción de costo: hoy el sistema corre sobre los niveles gratuitos
  de sus proveedores (base de datos, ruteo, clima). Escalarlo a un
  entorno productivo real tiene un costo mensual que el Capítulo 5 sí
  contempla dentro de la proyección económica, pero que el proyecto
  todavía no paga porque no genera ingresos reales.
\end{itemize}

\section{Matriz de riesgos}

Los riesgos con impacto económico directo ya están en la Tabla 5.9 del
Capítulo 5. Esta matriz reúne los demás: técnicos, de alcance, de
terceros, de datos y cumplimiento, y de producto (persuadir al usuario
de usar la aplicación y diferenciarse de cualquier alternativa),
ordenados de mayor a menor exposición.

\begin{footnotesize}
\begin{longtable}[]{@{}
  >{\raggedright\arraybackslash}p{(\columnwidth - 16\tabcolsep) * \real{0.0489}}
  >{\raggedright\arraybackslash}p{(\columnwidth - 16\tabcolsep) * \real{0.1007}}
  >{\raggedright\arraybackslash}p{(\columnwidth - 16\tabcolsep) * \real{0.3365}}
  >{\raggedright\arraybackslash}p{(\columnwidth - 16\tabcolsep) * \real{0.0992}}
  >{\raggedright\arraybackslash}p{(\columnwidth - 16\tabcolsep) * \real{0.0305}}
  >{\raggedright\arraybackslash}p{(\columnwidth - 16\tabcolsep) * \real{0.0304}}
  >{\raggedright\arraybackslash}p{(\columnwidth - 16\tabcolsep) * \real{0.0420}}
  >{\raggedright\arraybackslash}p{(\columnwidth - 16\tabcolsep) * \real{0.0791}}
  >{\raggedright\arraybackslash}p{(\columnwidth - 16\tabcolsep) * \real{0.2329}}@{}}
\caption{Matriz de riesgos de MotoSafe GBA, ordenada por exposición (probabilidad × impacto).}\label{tab:6-1}\\
\toprule\noalign{}
\begin{minipage}[b]{\linewidth}\raggedright
\textbf{ID}
\end{minipage} & \begin{minipage}[b]{\linewidth}\raggedright
\textbf{Categoría}
\end{minipage} & \begin{minipage}[b]{\linewidth}\raggedright
\textbf{Descripción (causa -- evento -- efecto)}
\end{minipage} & \begin{minipage}[b]{\linewidth}\raggedright
\textbf{Etapa}
\end{minipage} & \begin{minipage}[b]{\linewidth}\raggedright
\textbf{P}
\end{minipage} & \begin{minipage}[b]{\linewidth}\raggedright
\textbf{I}
\end{minipage} & \begin{minipage}[b]{\linewidth}\raggedright
\textbf{Exp.}
\end{minipage} & \begin{minipage}[b]{\linewidth}\raggedright
\textbf{Respuesta}
\end{minipage} & \begin{minipage}[b]{\linewidth}\raggedright
\textbf{Acción tomada o prevista}
\end{minipage} \\
\midrule\noalign{}
\endfirsthead
\toprule\noalign{}
\begin{minipage}[b]{\linewidth}\raggedright
\textbf{ID}
\end{minipage} & \begin{minipage}[b]{\linewidth}\raggedright
\textbf{Categoría}
\end{minipage} & \begin{minipage}[b]{\linewidth}\raggedright
\textbf{Descripción (causa -- evento -- efecto)}
\end{minipage} & \begin{minipage}[b]{\linewidth}\raggedright
\textbf{Etapa}
\end{minipage} & \begin{minipage}[b]{\linewidth}\raggedright
\textbf{P}
\end{minipage} & \begin{minipage}[b]{\linewidth}\raggedright
\textbf{I}
\end{minipage} & \begin{minipage}[b]{\linewidth}\raggedright
\textbf{Exp.}
\end{minipage} & \begin{minipage}[b]{\linewidth}\raggedright
\textbf{Respuesta}
\end{minipage} & \begin{minipage}[b]{\linewidth}\raggedright
\textbf{Acción tomada o prevista}
\end{minipage} \\
\midrule\noalign{}
\endhead
\bottomrule\noalign{}
\endlastfoot
R-01 & Datos y cumplimiento & Debido a que se sumó Google Noticias como
fuente de scraping sin revisar antes si el sitio lo permitía, se
incorporó una fuente que su propio sitio prohíbe usar de esa forma, y
que llegó a aportar la mayor parte del contenido nuevo de una corrida
real. & Durante el desarrollo & 5 & 4 & 20 & Mitigar (ya aplicada) & Ya
se retiró Google Noticias del scraper y se documentó el cambio. Como
prevención hacia adelante, revisar los términos de uso de cualquier
fuente nueva antes de sumarla. \\
R-02 & Alcance & Debido a que solo la Ciudad de Buenos Aires publica un
registro de siniestros con ubicación exacta (la Provincia de Buenos
Aires solo publica informes generales; existe una fuente nacional
parcial ---el Sistema de Alerta Temprana de Muertes Viales del SNIC,
geocodificada por partido--- pero limitada a víctimas fatales, sin
cubrir todas las severidades como el dataset de CABA), el modelo de
riesgo se entrenó únicamente con datos de CABA. Otras partes del sistema
sí cubren todo el GBA sin esa limitación (los reportes de la comunidad,
las notas de prensa georreferenciadas y el ruteo con tráfico funcionan
en cualquier punto con GPS), pero el cálculo de riesgo por zona, que es
el corazón del proyecto, todavía solo está validado para CABA, aunque el
nombre del proyecto prometa el GBA completo. & Antes de arrancar & 5 & 4
& 20 & Mitigar & Sostener en toda entrega que el modelo de riesgo está
validado solo para CABA, y que el resto del GBA se ofrece hoy como zona
de navegación, reportes y prensa, no de predicción. Ya se identificó una
fuente nacional parcial (SNIC, solo víctimas fatales, sin cubrir
heridos); evaluar si vale la pena integrarla queda pendiente de
decisión, no reemplaza al dataset de CABA. \\
R-03 & Datos y cumplimiento & Debido a que los primeros reportes
comunitarios se identificaban con un código que la propia aplicación
mostraba en público, sin ningún control de quién era el dueño del
reporte, cualquier usuario que viera ese código podía borrar un reporte
ajeno, lo que habría provocado pérdida de información de otras personas
sin ninguna forma de detectarlo. & Durante el desarrollo & 4 & 4 & 16 &
Mitigar (ya aplicada) & Se reemplazó el identificador público por un
token privado de una sola entrega para poder borrar un reporte. Los
reportes creados antes del cambio no tienen ese token y se mantienen
así, sin una migración retroactiva: es una limitación aceptada, no una
tarea pendiente. \\
R-04 & Producto y experiencia de usuario & Debido a que probar una
aplicación una sola vez no garantiza que el usuario vuelva a abrirla, y
a que la prueba real con los 12 motociclistas ya mostró fricción
concreta (comparar rutas fue la tarea que más ayuda necesitó, y la
demanda más repetida fue mostrar menos información por default), existe
el riesgo de que MotoSafe GBA no logre persuadir a los motociclistas de
incorporarlo a su rutina diaria, y quede como una aplicación que se
prueba una vez y se abandona, sin llegar a construir una base de
usuarios activa. & Operación & 4 & 4 & 16 & Mitigar & Aplicada. La
cobertura por fuente y el desglose de puntos de cada tarjeta de ruta
ahora quedan detrás de un control opcional ("Ver cómo se calcula"); por
default solo se muestra tiempo, distancia y la advertencia principal,
tal como pidieron los usuarios de la prueba real. Falta repetir la
validación con usuarios reales para confirmar que la fricción bajó. \\
R-05 & Producto y experiencia de usuario & Debido a que la aplicación
muestra un puntaje de riesgo por zona sin que el término alcance por sí
solo para despejar toda ambigüedad, en la prueba real con 12
motociclistas dos de ellos interpretaron ese puntaje como una predicción
de que a ellos les iba a pasar un accidente, en vez de una medida
relativa entre zonas. Si esa confusión se repite a mayor escala, los
usuarios podrían tomar decisiones basadas en una lectura equivocada del
dato, o perder confianza en la aplicación al notar la diferencia entre
lo que entendieron y lo que el sistema realmente mide. & Operación & 4 &
3 & 12 & Mitigar & Aplicada parcialmente. El texto aclaratorio dependía
de un tooltip de mouse, inútil en celular; se reemplazó por un ejemplo
visual siempre disponible con números reales ("cobertura de siniestros
40\% significa 4 de cada 10 tramos..."), visible junto al control de
detalle de la tarjeta de ruta. Falta repetir la prueba con usuarios
reales para confirmar si el ejemplo reduce la confusión detectada. \\
R-06 & De terceros & Debido a que el sistema depende de varios servicios
externos gratuitos sin garantía de disponibilidad (el clima, la
ubicación de direcciones y el ruteo con tráfico en vivo), si alguno de
ellos cae o cambia sus condiciones de uso, funciones centrales de la
aplicación (trazar una ruta, mostrar el clima de la zona, ubicar un
reporte por dirección) dejan de funcionar para cualquier usuario real,
no solo durante una demostración puntual. & Operación & 3 & 4 & 12 &
Mitigar & Sumar manejo de error visible cuando un servicio externo no
responde, en vez de dejar la función colgada sin aviso. Para la defensa
oral en particular, grabar de antemano un video corto de respaldo de
cada función principal por si algún servicio falla ese día. Si el
proyecto avanza a producción real, evaluar un nivel pago con garantía de
disponibilidad para las funciones más críticas. \\
R-07 & Producto y mercado & Debido a que MotoSafe GBA no es la única
forma posible de resolver el problema de información preventiva vial
(una aplicación de navegación ya instalada podría sumar una función
parecida, o podría aparecer un proyecto equivalente con más recursos de
desarrollo o de difusión), existe el riesgo de que, aunque el sistema
funcione bien técnicamente, no logre diferenciarse lo suficiente frente
a una alternativa más conocida, y quede sin la base de usuarios
necesaria para sostenerse más allá de la entrega académica. & Pasaje a
producción & 3 & 3 & 9 & Mitigar & Sostener como diferencial el foco
hiperlocal en motociclistas del GBA (reportes comunitarios, corredores
de riesgo, prensa georreferenciada combinados en un mismo mapa), algo
que hoy ninguna alternativa conocida ofrece junto; validar ese
diferencial con usuarios reales antes de pensar en escalarlo más allá de
la tesis. \\
R-08 & Técnicos & Debido a que el archivo de dependencias del proyecto
no fija una versión exacta de la librería de machine learning usada,
sino un rango abierto, una futura reinstalación del entorno podría traer
una versión más nueva no probada contra el modelo ya entrenado, lo que
podría hacer que las cifras de desempeño documentadas en el Capítulo 4
dejen de poder reproducirse exactamente tal como fueron medidas. &
Operación & 3 & 2 & 6 & Aceptar & Si en algún momento hace falta
reproducir el entrenamiento con exactitud (por ejemplo, ante una
revisión académica puntual), fijar en ese momento la versión exacta de
la librería. \\
\end{longtable}
\end{footnotesize}

\section{Riesgos de mayor exposición}

Los riesgos R-01 y R-02 comparten la mayor exposición de la matriz,
aunque por motivos distintos. R-01 es una falla que ya se produjo y ya
se corrigió: se sumó Google Noticias como fuente de scraping sin revisar
antes si el sitio lo permitía, y esa fuente llegó a aportar la mayor
parte del contenido nuevo de una corrida real antes de que se la
retirara. R-02 no es un error a corregir sino una limitación conocida
desde el inicio del desarrollo: no existe todavía, para el resto del
conurbano bonaerense, un registro de siniestros con la misma calidad que
el de CABA, así que el modelo de riesgo se entrenó solo con datos de esa
ciudad, aunque el nombre del proyecto prometa una cobertura mayor.
Nombrarlo como riesgo en este capítulo, y no solo mencionarlo al pasar,
sirve para no dar por sentada una cobertura que el sistema todavía no
puede sostener con evidencia.

R-01 y R-03 los dos son fallas que ya se produjeron y ya se corrigieron.
En ambos casos la corrección ya está aplicada y verificada; lo que se
agrega en este capítulo es el criterio para que el mismo tipo de error
no se repita: revisar los permisos de uso de una fuente antes de
sumarla, y no mostrar en público ningún código que funcione como llave
para borrar contenido.

Dos riesgos de este capítulo no son técnicos sino de producto, y se
tratan con el mismo peso que los anteriores porque tocan la razón de ser
del proyecto: que alguien lo use de verdad. R-04 recoge algo que ya se
observó en la prueba real con usuarios: probar la aplicación una vez no
garantiza que el motociclista vuelva a abrirla, y esa misma prueba
mostró fricción concreta en una de sus funciones centrales. R-07
reconoce que MotoSafe GBA compite, aunque sea de forma indirecta, contra
cualquier alternativa que un motociclista ya tenga instalada para
orientarse en la calle: el proyecto no gana usuarios solo por funcionar
bien por dentro, sino por ofrecer algo que esas alternativas no ofrecen.

\section{Riesgos estructurales del proyecto}

Además de la matriz anterior, conviene revisar el proyecto contra los
patrones generales que suelen hacer fracasar un proyecto de este tipo,
vistos en la cátedra, para dejar constancia de cuáles están presentes y
cuáles no:

\textbf{El alcance nunca termina de cerrar.} Presente en forma acotada:
el alcance creció varias veces durante el desarrollo (capa de
inseguridad, reportes de prensa, tráfico en vivo, corredores de riesgo),
pero cada expansión quedó documentada por etapa, no como un crecimiento
silencioso sin control.

\textbf{Depender de que alguien externo siga aportando datos.} No
aplica: los datos de siniestros son un registro público oficial
descargado una sola vez, sin depender de que una persona o entidad siga
entregando información para que el proyecto funcione.

\textbf{El equipo se desarma.} No aplica en el sentido tradicional,
porque el desarrollo está a cargo de una sola persona, pero sí existe un
riesgo parecido de coordinación: el proyecto se desarrolla con
asistencia de herramientas de inteligencia artificial, y no siempre se
revisó con el mismo nivel de detalle cada cambio sugerido antes de
incorporarlo, lo que contribuyó a las fallas de R-01 y R-03.

\textbf{La tecnología elegida no funciona como se esperaba.} No se
materializó: cada cambio de tecnología del proyecto (el motor de mapa,
el método de agrupamiento de zonas, el motor de rutas) fue una decisión
evaluada y documentada, no un cambio forzado porque algo dejó de andar.

\textbf{El costo o la infraestructura no son viables.} Presente y ya
tratado en el Capítulo 5: sostener el sistema en un entorno productivo
real depende de servicios pagos cuyo costo todavía no está validado
contra una fuente de ingresos real.

\textbf{Hay una barrera legal o ética de fondo. Presente de forma
acotada en R-01 (una fuente cuyo uso está restringido). La falta de una
política de privacidad formal para los datos de ubicación de los
reportes comunitarios se cerró con una política documentada: describe
qué se guarda (lat/lng exactas, categoría, autor\_id como UUID de
dispositivo, sin dato personal identificable), qué no se guarda (IP
persistida, nombre, email), retención y borrado por categoría, y declara
de forma explícita la limitación de que la ubicación se guarda sin
redondear. Esa política, escrita para los reportes comunitarios, todavía
no cubre los datos de cuenta que se sumaron después con el sistema de
autenticación (4.8): email, nombre de usuario, hash de contraseña
(bcrypt) y los códigos de verificación de un solo uso. Esos datos se
guardan en las tablas usuarios y codigos\_auth, se usan exclusivamente
para autenticar (nunca se muestran en una interfaz pública ni se
comparten con terceros), y falta extender el documento de política para
declararlos junto con los de reportes.}

\textbf{Nadie está a cargo del proyecto. No aplica: hay una sola persona
responsable de todas las decisiones, con registro de decisiones y de
cada paso en la bitácora de desarrollo.}

\section{Síntesis}

De los ocho riesgos relevados, dos ya se materializaron y fueron
tratados con evidencia verificable (R-01, R-03), lo que muestra que
MotoSafe GBA gestiona sus riesgos a medida que aparecen, en vez de
dejarlos para el informe final. De los seis restantes, dos siguen
abiertos con una acción concreta pendiente (R-02, R-07), dos tienen su
acción técnica aplicada y solo falta repetir la validación con usuarios
para confirmar el efecto (R-04, R-05, ya aplicadas), uno depende de
servicios externos y solo puede reducirse en parte (R-06), y uno se
acepta de forma explícita y justificada (R-08), porque el costo de
eliminarlo por completo hoy es mayor que su impacto real.

El riesgo de mayor exposición del proyecto no es de infraestructura ni
de plazos, sino de honestidad metodológica y de alcance real: haber
corregido a tiempo una fuente de datos que no estaba autorizada a usarse
(R-01), y sostener con claridad que el modelo de riesgo todavía solo
está validado para CABA, aunque el proyecto lleve GBA en el nombre
(R-02). Declarar estas dos cosas con evidencia, en vez de minimizarlas,
es en sí mismo parte de cómo se tratan: una tesis que reconoce sus
propios límites es más sólida frente al jurado que una que los esconde y
deja que otro los encuentre.

Pero un sistema puede estar técnicamente sólido y de todos modos no
conseguir usuarios reales. Los riesgos R-04 y R-07 quedan documentados
con ese mismo peso: el proyecto no se sostiene solo por funcionar bien
por dentro, sino por lograr que un motociclista real lo elija y lo siga
usando en vez de cualquier otra alternativa que ya tenga a mano.

\section{Seguimiento de los riesgos}

Esta matriz no se completa una sola vez y se archiva. El seguimiento se
hace en el propio repositorio del proyecto, donde cada hallazgo real
queda con su estado (validado, en curso, pendiente), en vez de vivir
solo en este capítulo. Algunos de los riesgos de esta matriz ya se
revisaron más de una vez a medida que apareció evidencia nueva durante
el desarrollo (el caso más claro es la medición de precisión del scraper
de prensa, corregida después de encontrar falsos positivos geográficos
reproducibles), y el mismo criterio se sostiene hasta la entrega final.
